\documentclass[a4paper,12pt]{article}
\usepackage[italian]{babel}
\usepackage[utf8]{inputenc}
\usepackage[T1]{fontenc}
\usepackage{geometry}
\usepackage{setspace}
\geometry{margin=2.5cm}
\onehalfspacing

\title{Verifica scritta di informatica - Classe 4 AIQ - Compito B - SOLUZIONI}
\author{Prof. Fedeli Massimo}
\date{}

\begin{document}
	\maketitle
	
	\section*{Soluzioni Parte A -- SQL}
	
	\begin{enumerate}
		
		\item
		\begin{verbatim}
			SELECT s.nome, s.cognome, c.titolo, i.data_iscrizione
			FROM STUDENTI s
			JOIN ISCRIZIONI i ON s.id_studente = i.id_studente
			JOIN CORSI c ON c.id_corso = i.id_corso
			WHERE s.id_studente = 5
			ORDER BY i.data_iscrizione DESC;
		\end{verbatim}
		
		\item
		\begin{verbatim}
			SELECT titolo
			FROM CORSI
			WHERE categoria = 'Informatica'
			AND attivo = 1
			AND costo > 100;
		\end{verbatim}
		
		\item
		\begin{verbatim}
			SELECT s.cognome, s.nome, c.titolo
			FROM STUDENTI s
			JOIN ISCRIZIONI i ON s.id_studente = i.id_studente
			JOIN CORSI c ON c.id_corso = i.id_corso
			WHERE c.livello = 'Base';
		\end{verbatim}
		
		\item
		\begin{verbatim}
			SELECT c.titolo
			FROM CORSI c
			JOIN ISCRIZIONI i ON c.id_corso = i.id_corso
			GROUP BY c.id_corso, c.titolo
			HAVING COUNT(*) >= 10;
		\end{verbatim}
		
		\item
		\begin{verbatim}
			SELECT i.id_iscrizione, s.cognome, i.data_iscrizione
			FROM ISCRIZIONI i
			JOIN STUDENTI s ON s.id_studente = i.id_studente
			WHERE i.id_iscrizione NOT IN (
			SELECT id_iscrizione
			FROM PAGAMENTI
			WHERE esito = 'positivo'
			);
		\end{verbatim}
		
		\item
		\begin{verbatim}
			SELECT c.titolo, COUNT(*) AS numero_iscritti
			FROM CORSI c
			JOIN ISCRIZIONI i ON c.id_corso = i.id_corso
			GROUP BY c.id_corso, c.titolo
			HAVING COUNT(*) >= 5;
		\end{verbatim}
		
		\item
		\begin{verbatim}
			SELECT s.cognome, s.nome, SUM(p.importo) AS spesa_totale
			FROM STUDENTI s
			JOIN ISCRIZIONI i ON s.id_studente = i.id_studente
			JOIN PAGAMENTI p ON p.id_iscrizione = i.id_iscrizione
			WHERE p.esito = 'positivo'
			GROUP BY s.id_studente, s.cognome, s.nome
			ORDER BY spesa_totale DESC;
		\end{verbatim}
		
		\item
		\begin{verbatim}
			SELECT c.titolo, AVG(r.voto) AS media_voti
			FROM CORSI c
			JOIN RECENSIONI r ON c.id_corso = r.id_corso
			GROUP BY c.id_corso, c.titolo;
		\end{verbatim}
		
		\item
		\begin{verbatim}
			SELECT s.id_studente, s.cognome
			FROM STUDENTI s
			JOIN ISCRIZIONI i ON s.id_studente = i.id_studente
			JOIN CORSI c ON c.id_corso = i.id_corso
			WHERE c.attivo = 1
			GROUP BY s.id_studente, s.cognome
			HAVING COUNT(DISTINCT c.id_corso) = (
			SELECT COUNT(*) FROM CORSI WHERE attivo = 1
			);
		\end{verbatim}
		
		\item
		\begin{verbatim}
			SELECT c.titolo, SUM(p.importo) AS incasso_totale
			FROM CORSI c
			JOIN ISCRIZIONI i ON c.id_corso = i.id_corso
			JOIN PAGAMENTI p ON p.id_iscrizione = i.id_iscrizione
			WHERE p.esito = 'positivo'
			GROUP BY c.id_corso, c.titolo;
		\end{verbatim}
		
	\end{enumerate}
	
	\section*{Soluzioni Parte B -- Algebra Relazionale}
	
	\begin{enumerate}
		
		\item
		\[
		\pi_{titolo}(\sigma_{attivo=1 \land livello='Avanzato'}(CORSI))
		\]
		
		\item
		\[
		\pi_{cognome, nome}(\sigma_{data\_iscrizione \geq '2026-01-01' \land data\_iscrizione \leq '2026-12-31'}(STUDENTI))
		\]
		
		\item
		\[
		\pi_{titolo}(CORSI \bowtie ISCRIZIONI)
		\]
		
		\item
		\[
		\pi_{cognome, nome, titolo}(
		\sigma_{data\_iscrizione \geq '2026-01-01' \land data\_iscrizione \leq '2026-01-31'}
		(STUDENTI \bowtie ISCRIZIONI \bowtie CORSI))
		\]
		
	\end{enumerate}
	
\end{document}